\documentclass[conference]{IEEEtran}
% The installed IEEEconf.cls is an unrelated IEEE Computer Society class.
% See README.md for template migration instructions.
\usepackage{amsmath,amssymb,amsthm}
\usepackage{graphicx}
\usepackage{cite}
\usepackage{booktabs}
\graphicspath{{figures/}}
\newcommand{\R}{\mathbb{R}}
\newcommand{\cN}{\mathcal{N}}
\newcommand{\norm}[1]{\lVert#1\rVert}
\newcommand{\avg}[1]{\overline{#1}}
\newtheorem{proposition}{Proposition}
\newtheorem{remark}{Remark}

\title{Self-Sustained Rotation in Cooperative Fencing\\with Average Velocity Tracking}
\author{\IEEEauthorblockN{Author Name(s)}
\IEEEauthorblockA{Affiliation and e-mail to be supplied before submission}}
\begin{document}
\maketitle
\begin{abstract}
An observer-based fencing controller for a moving target can achieve exponential convergence of the formation centroid while individual vehicles continue to move relative to the target. We examine this apparent discrepancy for the first controller of Kou, Chen, and Xiang. Exact decomposition of the closed loop into centroid and shape dynamics identifies the unregulated motion hidden from the average tracking objective. For every nondegenerate rigid rotation about the target, cancellation of the total torque of pairwise central repulsion forces selects the angular speed $|\omega|=\sqrt{k_2}$, independently of the formation shape. We also construct a regular-polygon rotating solution and characterize the neutral phase direction in a corotating linearization. Numerical results separate three long-time behaviours: rigid rotation with the predicted speed, the same shape frozen once the observer is modified, and nonrigid breathing with no common angular speed; a collinear branch instead closes onto the collision barrier at rate $1/t$. These observations do not establish global convergence to rotation; they show why centroid tracking alone cannot imply stationary relative positions.
\end{abstract}
\begin{IEEEkeywords}
Cooperative fencing, collective motion, rotating formations, observer dynamics, relative equilibria.
\end{IEEEkeywords}

\section{Introduction}
\label{sec:intro}
Cooperative target fencing asks a group of vehicles to keep a moving target in their convex hull while avoiding collisions. Unlike a prescribed rigid formation, the shape may be generated by local interactions. Kou, Chen, and Xiang~\cite{kou2022} introduced two controllers for a constant-velocity target whose velocity is unknown. Their first controller ensures convergence of the \emph{average} position and velocity errors. Its simulations also display sustained rotation about the target. Convergence of the mean is not convergence of each vehicle: a rotating configuration can have zero mean velocity error even though no individual vehicle stops in the target frame.

Circular collective motion is familiar in multiagent control. Sepulchre, Paley, and Leonard~\cite{sepulchre2007} analyzed circular relative equilibria of steered particles, while Ceccarelli \emph{et al.}~\cite{ceccarelli2008} designed circular motion around a reference beacon. Phase-weighted consensus can likewise prescribe a rotating formation~\cite{lin2010}. In a different setting, inconsistent interagent measurements can produce undesired collective motion~\cite{garcia2015}. In contrast, the rotation studied here needs neither a prescribed angular velocity nor a measurement mismatch: it is compatible with the nominal observer and repulsion law of~\cite{kou2022}. Other target-fencing formulations, such as double-integrator systems with distributed target observers~\cite{pan2024}, address different tracking and network requirements; their guarantees should not be conflated with individual stationarity for the present controller.

This paper makes three contributions. First, an exact centroid--shape decomposition clarifies what is and is not controlled by average target tracking. Second, a torque argument gives a necessary angular-speed law for \emph{any} nondegenerate rigid rotation, together with a constructive rotating branch. Third, a corotating linearization provides a precise, branch-dependent local stability test, complemented by numerical evidence of both rigid and nonrigid long-time motion. We do not claim that all collision-free trajectories become periodic or that rotation is globally attracting.

\section{Model and the Average-Tracking Gap}
\label{sec:model}
Consider $N\geq3$ identical planar single-integrator vehicles, with $\dot x_i=u_i$, $x_i\in\R^2$, and a target with unknown constant velocity $\dot x_0=v_0$. The first controller of~\cite{kou2022} is
\begin{align}
 u_i&=\phi_i+k_1(x_0-x_i)+v_i, &
 \dot v_i&=k_2(x_0-x_i),\label{eq:control}\\[-1.5ex]
 \phi_i&=\sum_{j\in\cN_i}\alpha(s_{ij})\frac{x_i-x_j}{s_{ij}}, &
 s_{ij}&=\norm{x_i-x_j},\label{eq:force}
\end{align}
where $k_1,k_2>0$, $\cN_i=\{j\ne i:s_{ij}\leq\mu\}$, and $d<\mu$ are the collision and sensing distances. The repulsion magnitude used in the numerical study is
\begin{equation}
 \alpha(s)=\begin{cases}
 (s-d)^{-1}-(\mu-d)^{-1},&d<s\leq\mu,\\
 0,&s>\mu.
 \end{cases}\label{eq:alpha}
\end{equation}
The analytical frequency result only requires symmetric, pairwise central forces: $\phi_i$ and $\phi_j$ must contain equal and opposite contributions parallel to $x_i-x_j$. Its derivation is valid along any collision-free rigid trajectory, including at a sensing cutoff as long as the force remains continuous.

Introduce target-frame variables $\xi_i=x_i-x_0$ and $\eta_i=v_i-v_0$. Since $v_0$ is constant,
\begin{equation}
 \dot\xi_i=\phi_i(\xi)-k_1\xi_i+\eta_i,
 \qquad \dot\eta_i=-k_2\xi_i.\label{eq:relative}
\end{equation}
For $\avg\xi=N^{-1}\sum_i\xi_i$ and $\avg\eta=N^{-1}\sum_i\eta_i$, pairwise antisymmetry yields $\sum_i\phi_i=0$ and hence
\begin{equation}
 \dot{\avg\xi}=-k_1\avg\xi+\avg\eta,
 \quad\dot{\avg\eta}=-k_2\avg\xi,
 \quad\ddot{\avg\xi}+k_1\dot{\avg\xi}+k_2\avg\xi=0.
 \label{eq:centroid}
\end{equation}
Both averages decay exponentially for positive gains. In particular, the target-to-convex-hull distance is no larger than $\norm{\avg\xi}$ because the mean vehicle position belongs to the convex hull. Also $N^{-1}\sum_i u_i\to v_0$. These facts say nothing about $u_i-v_0$ separately.

Set $q_i=\xi_i-\avg\xi$ and $z_i=\eta_i-\avg\eta$; then $\sum_iq_i=\sum_i z_i=0$. Translational invariance of $\phi$ implies that the shape subsystem is exactly
\begin{equation}
 \dot q_i=\phi_i(q)-k_1q_i+z_i,
 \qquad\dot z_i=-k_2q_i.\label{eq:shape}
\end{equation}
Thus the stable two-dimensional centroid oscillator is decoupled from $2(N-1)$ position-shape coordinates and their observer states. \emph{Decoupled does not mean undamped}: the fate of the shape variables depends on the interaction law. This distinction matters when explaining sustained motion.

\subsection{Two distinct notions of tracking}
The implication of~\eqref{eq:centroid} for the original convex-hull objective is especially simple. Since $x_0+\avg\xi=N^{-1}\sum_i x_i$ lies in $\operatorname{co}\{x_1,\ldots,x_N\}$, the target-to-hull distance is bounded by the norm of the stable oscillator state. This observation does not require a circular configuration, a fixed neighbor graph, or convergence of any individual $\xi_i$. Similarly, averaging the commanded inputs gives $\avg u=v_0+\dot{\avg\xi}\to v_0$. A stronger requirement, $u_i-v_0\to0$ for every $i$, is not a consequence of either average limit. On a nondegenerate rigid orbit, $u_i-v_0=\omega J\xi_i$ never decays even though $\avg u-v_0=0$ identically.

One can also see the distinction directly in~\eqref{eq:shape}. The average observer error is driven by $\avg\xi$, whereas every shape observer state $z_i$ integrates $-k_2q_i$. A periodic $q_i$ can generate a periodic $z_i$ with a quarter-cycle phase shift. The restoring term $-k_1q_i$ still acts in all directions; the shape dynamics must be analyzed before assigning any of their directions the label ``undamped.'' In particular, the mean equation cannot distinguish a rotating polygon from a nonrigid shape cycle with zero mean. Pairwise separations and individual angular speeds are needed to tell them apart.

\section{Rigid Rotations and Frequency Selection}
\label{sec:rotation}
Write $J=\left[\begin{smallmatrix}0&-1\\1&0\end{smallmatrix}\right]$ and $R(\theta)=\exp(J\theta)$. A rigid rotation about the target is a solution of~\eqref{eq:relative} of the form
\begin{equation}
 \xi_i(t)=R(\omega t)a_i,
 \quad \sum_i\norm{a_i}^2>0,\quad \omega\in\R,
 \label{eq:ansatz}
\end{equation}
with constant $a_i$. No common radius or regular polygon is assumed.

\begin{proposition}[Angular-speed selection]\label{prop:frequency}
Under symmetric pairwise central repulsion, every nondegenerate rigid rotation~\eqref{eq:ansatz} of~\eqref{eq:relative} has $\omega\ne0$ and
\begin{equation}
 \omega^2=k_2,\qquad
 \phi_i(a)=k_1a_i,\qquad \sum_i a_i=0.
 \label{eq:selection}
\end{equation}
Conversely, any collision-free nonzero configuration satisfying $\phi_i(a)=k_1a_i$ produces both signed rotations of~\eqref{eq:relative} if $\eta_i(0)=\omega Ja_i$ and $\omega=\pm\sqrt{k_2}$.
\end{proposition}
\begin{proof}
If $\omega=0$, all $\xi_i=a_i$ are constant, but $\eta_i(t)=\eta_i(0)-k_2ta_i$. The first equation of~\eqref{eq:relative} cannot then hold for a nonzero $a_i$ on a time interval. Therefore $\omega\ne0$. Rotational equivariance of the central forces gives $\phi_i(\xi(t))=R(\omega t)\phi_i(a)$. Integration of $\dot\eta_i=-k_2R(\omega t)a_i$ yields
\begin{equation}
 \eta_i(t)=(k_2/\omega)J\xi_i(t)+c_i.
 \label{eq:observerorbit}
\end{equation}
Inserting this into $\dot\xi_i=\omega J\xi_i$ shows that a fixed $c_i$ equals a rotating vector. Since $\omega\ne0$, it must vanish, giving
\begin{equation}
 (\omega-k_2/\omega)Ja_i=\phi_i(a)-k_1a_i.
 \label{eq:balance}
\end{equation}
Take the planar cross product of~\eqref{eq:balance} with $a_i$ and sum over $i$. Each pair contributes zero net torque about the target: $(a_i-a_j)\times[\alpha(s_{ij})(a_i-a_j)/s_{ij}]=0$. Consequently,
\begin{equation}
 (\omega-k_2/\omega)\sum_i\norm{a_i}^2
 =\sum_i a_i\times\phi_i(a)=0.
 \label{eq:torque}
\end{equation}
Nondegeneracy proves $\omega^2=k_2$, and~\eqref{eq:balance} reduces to radial balance. Summing the latter gives $k_1\sum_i a_i=\sum_i\phi_i(a)=0$. For the converse, define $\xi_i(t)=R(\omega t)a_i$ and $\eta_i(t)=\omega J\xi_i(t)$; both equations in~\eqref{eq:relative} hold because $\omega^2=k_2$.
\end{proof}

The condition on $\omega$ applies \emph{only} to rigid rotations. It does not impose a frequency on other periodic motions. It is a frequency-selection law, not an existence or global-attraction theorem. Since $k_2$ is a controller gain, the speed can be changed by changing $k_2$; it is not an independently specified maneuvering frequency. The sign is left free by reflection symmetry.

\subsection{An explicit rotating branch}
A concrete existence example follows from radial balance. For a centered regular hexagon of circumradius $r$ with $\mu/\sqrt3<r<\mu$ and $r>d$, only the two adjacent vertices interact. Their repulsive radial components add to $\phi_i=\alpha(r)a_i/r$; thus
\begin{equation}
 \alpha(r)=k_1r.\label{eq:hex}
\end{equation}
With $d=5$, $\mu=9$, and $k_1=0.5$, this equation has a root $r\simeq5.342$ in $(\mu/\sqrt3,\mu)$: $\alpha(\mu/\sqrt3)>k_1\mu/\sqrt3$, $\alpha(\mu)=0<k_1\mu$, and $\alpha$ decreases while $k_1r$ increases. Proposition~\ref{prop:frequency} therefore supplies an exact rigid orbit for either sign of $\omega$. This branch need not be attracting. In the baseline numerical run a different configuration, with alternating radii near $3.395$ and $6.312$, is selected. The angular-speed law survives this change of shape.

\subsection{Force balance versus frequency balance}
The proof separates two independent restrictions on the orbit. Internal central forces have zero resultant and zero total torque. Their zero torque eliminates the force term from~\eqref{eq:torque} and fixes $\omega^2=k_2$ without computing any interaction magnitude. Once this tangential condition holds, each individual equation~\eqref{eq:balance} requires $\phi_i(a)=k_1a_i$. The latter is a nonlinear algebraic shape equation: it depends on $N$, $d$, $\mu$, $k_1$, and the active neighbor graph. Thus even if two rotating branches have very different radii, their speeds must coincide at fixed $k_2$. Conversely, an arbitrary set of points with zero centroid generally fails radial balance and cannot be made into a closed-loop solution simply by choosing $\omega=\sqrt{k_2}$.

There is a useful identity for checking numerical equilibria. Taking the inner product of the radial-balance equations with $a_i$ and summing by unordered pairs gives
\begin{equation}
 k_1\sum_i\norm{a_i}^2
 =\sum_{\{i,j\}\,\mathrm{active}}\alpha(s_{ij})s_{ij}.
 \label{eq:radialidentity}
\end{equation}
Both sides are nonnegative, and the right side depends only on active separations. Identity~\eqref{eq:radialidentity} is necessary but not sufficient for a rotating shape: it checks the aggregate radial work, while~\eqref{eq:selection} checks force balance vehicle by vehicle. In computations we therefore inspect the residuals $\norm{\phi_i(a)-k_1a_i}$ and the drift of the pairwise distances in addition to the frequency fit.

The constructed hexagon is not a uniqueness result. Repulsion is radial for a regular polygon by symmetry, but the active edges change when a diagonal crosses $\mu$. A root of the algebraic balance equation must therefore be checked against the neighbor geometry used to derive it. For the parameters above, the next-nearest-neighbor separation is $\sqrt3r\simeq9.25>\mu$, so the two-neighbor formula is self-consistent. A numerical orbit selected from another initial condition may have a different active graph and a different distribution of radii. Its frequency can be predicted before its shape is known.

\section{Shape Energy and Corotating Stability}
\label{sec:stability}
The force is locally the gradient of a rotationally invariant potential $P(q)=\sum_{i<j}\Phi(s_{ij})$ with $\Phi'(s)=\alpha(s)$, away from collisions and sensing cutoffs. Eliminating $z$ from~\eqref{eq:shape} gives
\begin{equation}
 \ddot q+B(q)\dot q+k_2q=0,
 \quad B(q)=k_1I-\nabla^2P(q).
 \label{eq:secondorder}
\end{equation}
All matrices in this section act on the zero-mean shape space. The oscillator energy obeys
\begin{equation}
 E=\tfrac12\norm{\dot q}^2+\tfrac{k_2}{2}\norm q^2,
 \qquad \dot E=-\dot q^{\mathsf T}B(q)\dot q.
 \label{eq:energy}
\end{equation}
For an active edge of length $s$, $\nabla^2\Phi$ has radial and tangential curvatures $\alpha'(s)<0$ and $\alpha(s)/s>0$. Thus $B$ need not be positive semidefinite away from a rotating branch;~\eqref{eq:energy} is not a global Lyapunov certificate. At a radial-balance configuration $a$, differentiating $\phi(R(\theta)a)=R(\theta)\phi(a)$ yields $\nabla^2P(a)Ja=J\phi(a)=k_1Ja$, so
\begin{equation}
 B(a)Ja=0.\label{eq:shapezero}
\end{equation}
This identifies a symmetry-induced neutral direction, not an absence of damping in every shape direction.

For local analysis choose $\omega_s=\pm\sqrt{k_2}$ and let $Q=R(-\omega_st)q$ and $Z=R(-\omega_st)z$. On the zero-mean space the corotating dynamics are autonomous:
\begin{align}
 \dot Q&=\phi(Q)-k_1Q-\omega_s J_N Q+Z,\nonumber\\
 \dot Z&=-k_2Q-\omega_s J_N Z,
 \label{eq:corot}
\end{align}
where $J_N=I_N\otimes J$ restricted to zero-mean vectors. A radial-balance configuration $a$ is an equilibrium together with $Z^*=\omega_sJ_Na$. Define $H=\nabla^2P(a)-k_1I$. The linearization is
\begin{equation}
 A_s=\begin{bmatrix}H-\omega_sJ_N&I\\
 -\omega_s^2I&-\omega_sJ_N\end{bmatrix}.
 \label{eq:jacobian}
\end{equation}
Rotational invariance gives the exact zero eigenvector
\begin{equation}
 A_s\begin{bmatrix}J_Na\\-\omega_sa\end{bmatrix}=0.
 \label{eq:phase}
\end{equation}
If this zero eigenvalue is simple and every other eigenvalue of the \emph{full shape-space} linearization has strictly negative real part, the circle of equilibria in the corotating frame is locally exponentially attracting transversely, hence the rotating orbit is locally orbitally asymptotically stable. The stable centroid subsystem~\eqref{eq:centroid} does not alter this test. Neither~\eqref{eq:shapezero} alone nor the spectrum of a different branch is sufficient for that conclusion.

\subsection{Interpreting the phase mode}
The neutral direction in~\eqref{eq:phase} is not an accidental numerical eigenvalue. Rotating $a$ through a constant angle produces a continuous family of equilibria of~\eqref{eq:corot}; differentiation with respect to that angle yields $(J_Na,-\omega_sa)$. Perturbations along this direction change only the phase of the same periodic orbit. Consequently, ordinary asymptotic stability of an individual point in the corotating frame is impossible: the correct local notion is convergence to the \emph{orbit}, allowing a phase offset. Positive real parts transverse to the phase direction would indicate instability of the selected shape, while additional eigenvalues on the imaginary axis would make the simple orbital-attraction test inconclusive.

The test must be performed on the zero-mean subspace. Including mean coordinates would duplicate stable centroid modes and could obscure which zero eigenvalue is symmetry-induced. It must also use the Hessian at the \emph{selected} radial-balance configuration, not at a visually similar regular polygon. The repulsion law~\eqref{eq:alpha} is continuous but changes smooth formula at $s=\mu$; classical linearization assumes no equilibrium edge lies exactly at that cutoff. In the reported baseline branch the relevant edge lengths are away from the switching threshold. For an equilibrium at the cutoff one would instead need a nonsmooth or piecewise-smooth analysis.

Equation~\eqref{eq:energy} offers a complementary physical reading. At radial balance, covariance guarantees $B(a)J_Na=0$ and the rigid tangent velocity $\dot q=\omega J_Nq$ does not dissipate the oscillator energy. Away from that configuration, some tangential edge curvatures can exceed $k_1$ so that $B$ is indefinite. During a transient, the oscillator energy can then increase despite the positive direct gain $k_1$. This explains why a positive gain in the individual position feedback cannot, by itself, rule out self-sustained rotation. It does not, however, prove that an arbitrary initial shape reaches a particular rotating solution.

\section{Numerical Study}
\label{sec:sim}
The numerical evidence uses three complementary observables: the centroid trace, which tests the linear dynamics~\eqref{eq:centroid}; the minimum separation, which tests whether a trajectory stays on the collision-free side of $d$; and the complete set of pairwise separations, which tests rigidity. For each vehicle whose radius stays away from zero we also use the instantaneous angular speed
\begin{equation}
 \widehat\omega_i(t)=\frac{\xi_i(t)\times\dot\xi_i(t)}{\norm{\xi_i(t)}^2}.
 \label{eq:diagnostic}
\end{equation}
A rigid orbit requires stable pairwise distances \emph{and} a common $\widehat\omega_i$; a common Fourier peak of one single pairwise distance establishes neither.

Figure~\ref{fig:regimes} compares the first controller with the second one on the same gains and the same hexagon initial condition, and adds a run in which neither rigid motion nor a fixed shape survives. Each row is one run, each column a snapshot inside the asymptotic window, and the marker colour is the binned $\widehat\omega_i$.

\begin{figure*}[!t]
 \centering\includegraphics[width=.95\textwidth]{fig-regimes.pdf}
 \caption{Long-time behaviour in the target frame ($\xi_i=x_i-x_0$). Rows: rigid rotation (first controller), static formation (second controller of Sec.~\ref{sec:second}, same gains and the same hexagon initial condition), breathing (first controller, Case C). Columns: three snapshots inside the asymptotic window, with a $0.4$~s motion trail; the star is the target. Marker colour is the binned instantaneous angular speed~\eqref{eq:diagnostic}, not a vehicle identity. Common parameters: $N=6$ (rows 1--2), $d=5$, $\mu=9$, $k_1=k_2=0.5$, $\alpha$ from~\eqref{eq:alpha}, $v_0=(1,0)$, $\eta_i(0)=0$, LSODA at relative tolerance $10^{-9}$; row 3 uses $N=5$ and the case-C initial condition $\xi_i(0)=(-10.58,-12.53)$, $(-5.38,5.90)$, $(-0.18,9.18)$, $(-7.35,-4.81)$, $(-6.29,12.44)$ with $\eta_i(0)=(3.53,-1.27)$, $(-0.51,-1.49)$, $(1.97,-3.68)$, $(-3.46,-0.77)$, $(-2.04,2.76)$, integrated to relative tolerance $10^{-8}$.}
 \label{fig:regimes}
\end{figure*}

\emph{Rigid rotation.} The late angular speed is $-0.707107$, i.e.\ $-\sqrt{k_2}$, with a spread of $5\times10^{-8}$ between vehicles, so the diagnostic cannot separate them. The distinct final separations are $5.471$, $5.881$, $9.707$ and $10.932$, all above $d$, and the radii alternate between $3.395$ and $6.312$ rather than the regular-hexagon value $r^*\simeq5.342$ of~\eqref{eq:hex}. Sweeping $k_2=0.05,0.1,0.2,0.5,1,2$, a factor $40$ in $|\omega|$, keeps the relative discrepancy $|\widehat\omega_i|/\sqrt{k_2}-1$ below $2.4\times10^{-5}$, largest at the smallest gain where the transient is longest; this checks Proposition~\ref{prop:frequency} numerically and is not a convergence proof. For the same branch at $k_2=0.5$ the projected $20\times20$ matrix~\eqref{eq:jacobian} has one eigenvalue of magnitude below $3\times10^{-6}$ aligned with~\eqref{eq:phase}, and every other computed eigenvalue has negative real part with abscissa $-0.0780$; these are branch-specific numbers, not a certified spectral bound.

\emph{Static formation.} Replacing only the observer, as in Sec.~\ref{sec:second}, gives no rotation from the same initial condition and the same gains: all three snapshots coincide, $\widehat\omega_i=0$ (spread $2\times10^{-11}$), the observer error decays to $\max_i\norm{\eta_i}=4\times10^{-12}$, and the frozen shape is the \emph{same} unequal-radius configuration, with radii $3.395$ and $6.312$ and nearest separations $5.471$. The second controller therefore does not select a different shape; it removes the circulation while the shape is left intact, which is consistent with $\dot V=-\sum_i\norm{F_i}^2$ in~\eqref{eq:contrast}.

\emph{Breathing.} The case-C run keeps oscillating: pairwise separations have a dominant finite-window period of about $2.48$~s and, at a single instant, $\widehat\omega_i$ ranges over $[-2.45,0.12]$~rad/s, so no common angular speed exists and Proposition~\ref{prop:frequency} does not apply. Holding everything else fixed and raising $k_1$ to $0.7$ suppresses the oscillation and restores $\widehat\omega_i=-\sqrt{k_2}=-0.7071$ with zero spread; lowering $k_1$ to $0.3$ enlarges the excursions and brings one vehicle to a radius of $0.6$. These finite-horizon runs suggest a shape-gain-dependent boundary between breathing and rigid rotation; they neither prove a periodic attractor nor classify all trajectories.

\begin{figure}[!t]
 \centering\includegraphics[width=\columnwidth]{fig-jamrate.pdf}
 \caption{Jamming on the collision barrier, first controller, restricted to the collinear invariant subspace. The minimum separation minus the safety distance decreases with a fitted log-log slope of $-0.99$ and follows $c^*/t$ with $c^*=0.089$, while the observer state grows linearly; every $\widehat\omega_i$ vanishes. Planar non-collinear initial conditions did not reach this branch in our runs, so no jammed attractor in the plane is claimed.}
 \label{fig:jamrate}
\end{figure}

\emph{Jamming.} Figure~\ref{fig:jamrate} isolates a third mechanism of the first controller. On the collinear invariant subspace the shape converges while the minimum separation closes onto $d$ at rate $\min_{i<j}s_{ij}-d\simeq c^*/t$, here with $c^*=0.089$ and a fitted exponent $-0.99$, and the observer states diverge linearly. This branch violates the premise of Proposition~\ref{prop:frequency} in the strongest way: the formation does not rotate at all, yet the target is still fenced in the centroid sense. Because the branch lives on an invariant subspace reached only from collinear (or numerically collinear) data, we report it as a limiting case rather than as a planar attractor.

Two caveats apply throughout. The angular-speed diagnostic~\eqref{eq:diagnostic} is only meaningful while $\norm{\xi_i}$ stays away from zero, which is why the jammed run is quantified by its gaps. Since the solver works in the target frame, the constant target velocity affects only the lab-frame translation. The simulations are reproducible from the bundled script, which regenerates both figures and prints the numbers quoted above.

\section{Implications and Limitations}
\label{sec:discussion}
\subsection{Removing the rotating solution}
\label{sec:second}
The second controller of~\cite{kou2022} retains $u_i$ from~\eqref{eq:control} but changes the observer to $\dot v_i=\phi_i-k_1\xi_i$. For its target-frame dynamics define $F_i=\phi_i-k_1\xi_i$. With the repulsion-barrier potential $U=\sum_{i<j}\int_{s_{ij}}^{\mu}\alpha(s)\,ds+(k_1/2)\sum_i\norm{\xi_i}^2$, one has $\nabla_{\xi_i}U=-F_i$. The Lyapunov function $V=U+(1/2)\sum_i\norm{\eta_i}^2$ satisfies
\begin{equation}
 \dot V=-\sum_i\norm{F_i}^2.\label{eq:contrast}
\end{equation}
Here $\dot\xi_i=F_i+\eta_i$ and $\dot\eta_i=F_i$, so~\eqref{eq:contrast} follows by exact cancellation of cross terms. The first controller lacks this cancellation. Equation~\eqref{eq:contrast} excludes a nonstationary periodic relative orbit with $F_i\ne0$; on an invariant subset with $F_i=0$, maintaining $\dot\eta_i=0$ while $\dot\xi_i=\eta_i$ in a bounded collision-free configuration requires stationary relative motion. Full individual convergence relies on the regularity, boundedness and invariance assumptions of the original controller analysis~\cite{kou2022}, not merely the sign of $\dot V$.

\subsection{What remains unresolved}
The exact law~\eqref{eq:selection} is conditional on rigid motion. The existence construction shows one orbit, but its basin and the stability of other radial-balance shapes are distinct questions. In particular, tangential curvature can make $B(q)$ indefinite during transients, preventing~\eqref{eq:energy} from establishing global attraction. A signed circulation $L=\sum_i\xi_i\times\eta_i$ satisfies $L=\omega\sum_i\norm{\xi_i}^2$ \emph{on} a rigid orbit; away from it $L$ is not conserved, so the sign of the initial circulation is not a general direction-selection theorem. Finally, distance oscillations and boundary approach under special initial conditions require separate analytical treatment. We make no global boundedness or trichotomy claim for the first controller.

\section{Comparison With Related Motion Mechanisms}
\label{sec:comparison}
The literature on rotating formations spans several different dynamical questions. In deliberate circular-motion control, a reference center and a desired direction or turning rate are part of the objective. Ceccarelli \emph{et al.}~\cite{ceccarelli2008} seek circular motion around a virtual beacon for nonholonomic vehicles and study locally stable multi-vehicle equilibrium configurations. Sepulchre \emph{et al.}~\cite{sepulchre2007} regard parallel and circular collective motions as relative equilibria of unit-speed steered particles. Their geometric viewpoint is useful here: a circle of equilibria in rotating coordinates corresponds to one periodic orbit modulo phase in fixed coordinates. The present problem, however, is not to stabilize a prescribed turn. The control goal in~\cite{kou2022} is target fencing, and the angular-speed law emerges from the closed-loop observer dynamics after a rigid orbit is assumed.

A second distinction concerns what drives persistent motion. Phase-weighted formation control~\cite{lin2010} explicitly introduces the desired rotation into the protocol; it coordinates agents to execute that motion. Inconsistent distance measurements~\cite{garcia2015} can instead upset a nominally stationary rigid formation, leading to residual collective motion. Neither construction is required here: every agent can use the same correct target position and the same symmetric repulsion law. The remaining circulation is consistent with the objective precisely because only the mean velocity is required to track $v_0$. This is why a measurement calibration that removes mismatches would not address the observer--repulsion mechanism studied in this paper.

A third distinction lies in the word ``fencing.'' A convex-hull condition describes where the target sits relative to the vehicles, not how the vehicles move while meeting that condition. Work on target fencing with double-integrator vehicles and distributed target observers~\cite{pan2024} extends the task to other dynamics and information patterns. Such contributions should be compared on their explicit objective sets: a centroid or hull requirement, prescribed spacing, convergence of each velocity, and stationarity in the target frame are different statements. In the present first-controller model, the first two equations in~\eqref{eq:centroid} certify a mean property, while Proposition~\ref{prop:frequency} exhibits nonstationary solutions compatible with it.

\section{Additional Analytical Checks}
\label{sec:checks}
\subsection{Regular configurations and active-edge geometry}
The hexagon example in Sec.~\ref{sec:rotation} is a particularly transparent case because its chord lengths are $r$, $\sqrt3r$, and $2r$. For $\mu/\sqrt3<r<\mu$, only nearest neighbors contribute. The repulsive force from the two adjacent vertices is collinear with $a_i$ and has magnitude $\alpha(r)$, so the scalar balance~\eqref{eq:hex} is necessary and sufficient for that symmetric configuration. Since $\alpha$ is strictly decreasing on $(d,\mu)$ and $k_1r$ increases strictly, at most one root occurs within this active-edge regime. The endpoint inequalities stated after~\eqref{eq:hex} guarantee one root in this particular example; they are not automatic for all gains or sensing radii.

A change of $N$ changes both the polygon geometry and its active edges. For a regular $N$-gon, the nearest-neighbor chord is $s=2r\sin(\pi/N)$ and each nearest-neighbor force has radial projection $\alpha(s)\sin(\pi/N)$. Provided all other chords exceed $\mu$ and $d<s\leq\mu$, the force balance becomes
\begin{equation}
 2\alpha(s)\sin(\pi/N)=k_1r,
 \quad\text{or}\quad
 \alpha(s)=\frac{k_1s}{4\sin^2(\pi/N)}.
 \label{eq:polygon}
\end{equation}
This formula supplies candidate regular rotating branches by solving a one-dimensional equation, \emph{then} checking the assumed active-edge inequalities. If additional chords are sensed, their radial contributions must be added before using the condition. The speed result of Proposition~\ref{prop:frequency} is unaffected, because its proof sums torques over all active pairs without requiring nearest-neighbor interactions.

The observer initial condition in the converse part of Proposition~\ref{prop:frequency} is also essential. A radial-balance shape $a$ together with $\eta_i(0)=0$ is generally not initialized on its exact rotating orbit. The exact orbit requires $\eta_i(0)=\omega Ja_i$, corresponding to a tangential observer offset. A simulation begun with zero estimates must first select and approach a rotating branch, which is an attraction question, not an algebraic existence result. This distinction explains why numerical transients may depart substantially from a chosen regular polygon even when an exact regular solution exists at the same gain values.

\subsection{A reproducible Hessian assembly}
For a pair $(i,j)$ at separation vector $y=a_i-a_j$ and $s=\norm y$ in the active range, let $e=y/s$. Differentiating the pair force gives the symmetric $2\times2$ block
\begin{equation}
 K_{ij}=\alpha'(s)ee^{\mathsf T}
 +\frac{\alpha(s)}{s}(I-ee^{\mathsf T}).
 \label{eq:edgehessian}
\end{equation}
In the full Hessian $\nabla^2P(a)$ this block contributes $+K_{ij}$ to diagonal blocks $(i,i)$ and $(j,j)$, and $-K_{ij}$ to off-diagonal blocks $(i,j)$ and $(j,i)$. Inactive pairs contribute zero. The incidence-like block pattern automatically annihilates uniform translations. On the zero-mean subspace, choose a matrix $T$ whose columns form an orthonormal basis, so $T^{\mathsf T}T=I$ and $T^{\mathsf T}(\mathbf1_N\otimes b)=0$ for every planar $b$. The reduced matrices used in~\eqref{eq:jacobian} are
\begin{equation}
 H_r=T^{\mathsf T}(\nabla^2P(a)-k_1I)T,
 \qquad J_r=T^{\mathsf T}J_NT.
 \label{eq:project}
\end{equation}
Then $A_s$ is assembled from $H_r$, $J_r$, and identity blocks of size $2(N-1)$. For $N=6$ its order is $20$, as stated in Sec.~\ref{sec:sim}. Projecting is not a cosmetic step: an unprojected spectrum contains additional centroid eigenvalues and does not directly test stability of shape perturbations. An independent numerical check is to evaluate $\norm{H_rT^{\mathsf T}J_Na}$ near the candidate equilibrium and compare its phase vector with the computed eigenvector nearest zero.

\subsection{The limits of energy arguments}
At a stationary shape $a$ with $\phi(a)=k_1a$, the rotational symmetry identity~\eqref{eq:shapezero} implies a zero eigenvalue of the symmetric matrix $B(a)$. It is tempting to infer that $B(q)\succeq0$ in a neighborhood and then apply~\eqref{eq:energy} as a Lyapunov function. That inference is invalid: positive semidefiniteness at a point with a zero eigenvalue is not an open property. In addition, the tangential curvature $\alpha(s)/s$ increases sharply when a pair approaches the repulsion barrier. Near such a configuration, the indefinite part of $B(q)$ can inject energy into shape motion. The stability test in Sec.~\ref{sec:stability} therefore uses the \emph{complete corotating Jacobian}, not the sign of $B(a)$ alone. Establishing a larger basin would require additional estimates that control both the barrier geometry and changes in the active graph.

\section{Conclusion}
A pairwise-central-force fencing controller can track a target in the centroid sense yet retain nonzero individual motion. For every nondegenerate rigid rotation of the nominal closed loop, zero total internal torque and the integral observer enforce $|\omega|=\sqrt{k_2}$ and radial force balance; the angular-speed law does not depend on the particular rotating shape. A regular-hexagon branch establishes existence for the reported parameters, while corotating spectra and simulations indicate local attraction of a different, unequal-radius branch. Replacing only the observer freezes that same shape; other initial conditions oscillate without a rigid shape, and a collinear branch closes onto the collision barrier at rate $1/t$. Understanding attraction basins and a complete long-term classification remains open.

\begin{thebibliography}{9}
\bibitem{kou2022} L. Kou, Z. Chen, and J. Xiang, ``Cooperative fencing control of multiple vehicles for a moving target with an unknown velocity,'' \emph{IEEE Trans. Autom. Control}, vol. 67, no. 2, pp. 1008--1015, 2022, doi: 10.1109/TAC.2021.3075320.
\bibitem{sepulchre2007} R. Sepulchre, D. A. Paley, and N. E. Leonard, ``Stabilization of planar collective motion: All-to-all communication,'' \emph{IEEE Trans. Autom. Control}, 2007, doi: 10.1109/TAC.2007.898077.
\bibitem{ceccarelli2008} N. Ceccarelli, M. Di Marco, A. Garulli, and A. Giannitrapani, ``Collective circular motion of multi-vehicle systems,'' \emph{Automatica}, 2008, doi: 10.1016/j.automatica.2008.04.024.
\bibitem{lin2010} P. Lin and Y. Jia, ``Distributed rotating formation control of multi-agent systems,'' \emph{Syst. Control Lett.}, vol. 59, no. 10, pp. 587--595, 2010, doi: 10.1016/j.sysconle.2010.06.015.
\bibitem{garcia2015} H. Garcia de Marina, M. Cao, and B. Jayawardhana, ``Controlling rigid formations of mobile agents under inconsistent measurements,'' \emph{IEEE Trans. Robot.}, vol. 31, no. 1, pp. 31--39, 2015, doi: 10.1109/TRO.2014.2373145.
\bibitem{pan2024} Z. Pan and B. M. Chen, ``Cooperative target fencing of multiple double-integrator systems with connectivity preservation,'' \emph{Int. J. Robust Nonlinear Control}, 2024, doi: 10.1002/rnc.7381.
\end{thebibliography}
\end{document}
